%!TEX program = lualatex
\documentclass[ngerman,headheight=70pt]{scrartcl}
\usepackage[ngerman]{babel}
\usepackage{blindtext}
\usepackage[utf8]{luainputenc}
\usepackage{scrlayer-scrpage}
\usepackage[T1]{fontenc}
\usepackage{fontspec}
\usepackage{multicol}
\usepackage{enumitem}
\usepackage{epigraph}
\usepackage[german=quotes]{csquotes}
\usepackage[bottom=2cm,footskip=8mm]{geometry}
\setmainfont{Source Sans Pro}
\MakeOuterQuote{"}

\setlength{\parindent}{0pt}
\setlength{\parskip}{6pt}

\pagestyle{plain.scrheadings}
\thispagestyle{scrheadings}
\lohead{\vspace*{-1.5cm}\includegraphics[scale=0.4]{../images/stupa_logo_rot.pdf}}
\rohead{\\Präsidium des\\Studierendenparlaments}

\lofoot[Universität Hamburg · Präsidium des Studierendenparlaments\\\hrule Von-Melle-Park 5 · D-20146 Hamburg]{Universität Hamburg · Präsidium des Studierendenparlaments\\\hrule Von-Melle-Park 5 · D-20146 Hamburg}
\cofoot[]{}
\cohead[\thepage]{}

\renewcommand{\thesubsection}{TOP \arabic{subsection}}
\renewcommand{\thesubsubsection}{\arabic{subsubsection}.}
\setcounter{subsection}{-1}

\begin{document}
    UHH · StuPa-Präsidium · Von-Melle-Park 5 · D-20146 Hamburg

    \section*{Beschluss des Studierendenparlamentes vom 12. Mai 2016}
    \subsection*{Herren-Fußball-Europameisterschaft 2016 – und wieder
    Großsportveranstaltung}

    Spätestens seit der Auseinandersetzung um Olympische Spiele liegt für die
    meisten Hamburger offen, dass derartige kommerzielle Großevents vor allem
    ein großes Geschäft für Veranstalter, Sponsoren, Ausstatter und Wettbüros
    sind, auf Kosten der großen Mehrheit der Bevölkerung – kulturell begleitet
    von Leistungsideologie, Doping und nationalen Gefühlen.

    Und Medienhype: Keine Zeitung, die nicht mit Torjubel titelt, keine
    Nachrichten, ohne die neusten Informationen aus den Umkleidekabinen. Die
    gnadenlose Suggestion lautet: hältst Du Fußball (oder Handball, oder Laufen,
    oder Schießen, oder was auch immer) nicht für das wichtigste der Welt, bist
    Du ein durchgeknallter Freak. Ein Brief einer offensichtlich konservativen
    Leserin an die FAZ vom 17. Juni 2006 fasste das so zusammen:

    \blockquote[Renate Jansen, Düsseldorf]{Bitte, reduzieren Sie die Berichterstattung über diese duselige WM, ansonsten
    überlege ich mir ernsthaft, für die Zeit der WM eine F.A.Z.-Pause einzulegen.
    Mich interessieren einzelne oder ganze Körperteile von Spielern nicht,
    ebensowenig die Tipergebnisse von Ministern (haben die sonst nichts zu tun?)
    Bieten Sie bitte seriösen Journalismus, decken Sie die Schwächen der Regierung
    egal welcher Farbe auf und machen Sie vernünftige Vorschläge, wie die Karre
    Deutschland wieder aus dem Dreck kommt.}

    Der AStA der Uni Hamburg wird sich an dieser Gehirnwäsche nicht beteiligen,
    mit der den Menschen vorgeschrieben werden soll, was sie zu denken haben:
    dass ihnen ein Sportereignis das wichtigste im Leben sein möge.

    Um eine kulturelle Alternative zu schaffen, organisiert der AStA kurzfristig
    eine Veranstaltungsreihe. Angestrebt wird, rund um die Spieltage der
    Mannschaft aus Deutschland, jeweils einen Referenten aus dem Land der
    gegnerischen Mannschaft zu gewinnen, der von den dortigen sozialen
    Auseinandersetzungen berichtet.

    Darüber hinaus soll der AStA sich weiterhin bemühen, den Breitensport zu
    fördern und in diesem Zusammenhang die Durchführung eines
    AStA-Breitensport-Fussballturniers prüfen.
\end{document}
